% ATENÇÃO - veja com o seu orientador se você vai ter este capítulo e se este vai ter nome!
\chapter{Trabalhos Relacionados}
\label{cap:trabalhos:relacionados}

Este capítulo apresenta estudos acadêmicos e técnicos anteriores que realizaram análises comparativas ou avaliaram ferramentas de CI/CD. A discussão tem como objetivo destacar resultados e critérios considerados relevantes por diferentes autores sobre as ferramentas estudadas, contextualizando GitHub Actions, GitLab CI/CD e Jenkins.

\cite{Shahin2017} realizaram uma revisão sistemática sobre abordagens, ferramentas e desafios relacionados à CI. O estudo destaca que Jenkins e GitHub Actions são amplamente adotados devido à sua flexibilidade e à capacidade de integração com outros serviços. Complementarmente, \cite{RiunguKalliosaari2016} destacam que a adoção dessas ferramentas contribui significativamente para a redução de erros manuais e aprimoramento da colaboração entre as equipes.

Na análise comparativa feita por Graciano em \cite{GRACIANOKAKINOHANA2023} entre GitHub Actions e GitLab CI/CD, critérios como tempo de execução dos pipelines, facilidade de configuração e integração com outras ferramentas são abordados. Nesse estudo, GitHub Actions destaca-se pela integração fluida com o ecossistema GitHub e facilidade de uso, enquanto o GitLab CI/CD apresenta vantagens na flexibilidade da configuração e no suporte a runners dedicados, que permitem execução em infraestruturas próprias.

Segundo Lima em \cite{Lima2021}, ao comparar especificamente as funcionalidades de Jenkins e GitLab CI/CD, Jenkins oferece maior flexibilidade devido à sua natureza independente e extensa disponibilidade de plugins. No entanto, ressalta que isso resulta em uma complexidade maior na configuração e manutenção. Já o GitLab CI/CD é descrito como oferecendo uma experiência mais simplificada e intuitiva, particularmente para equipes que já utilizam o GitLab para gestão de código.

Estudos clássicos como o de \cite{Duvall2007} reforçam a importância de considerar fatores como tempo de execução dos pipelines, facilidade de uso e suporte abrangente para testes automatizados ao selecionar uma ferramenta de CI/CD. Adicionalmente, o relatório \textit{Best CI/CD Tools} \cite{bluelightCo} classifica GitHub Actions, GitLab CI/CD e Jenkins como líderes de mercado, destacando os pontos fortes e limitações de cada um e enfatizando que a escolha depende do nível de controle desejado sobre a infraestrutura e das características do ecossistema de desenvolvimento utilizado.

Esses estudos fornecem uma base importante que justifica a relevância da análise comparativa proposta neste trabalho. Ao considerar as funcionalidades, limitações, facilidade de configuração e o tempo exigido para implementação das pipelines, os resultados deste estudo buscam contribuir com orientações práticas para equipes de desenvolvimento na escolha da solução mais adequada às suas necessidades específicas.
